\slajdid{10-s066}
\begin{frame}[label=10-s066]{Dinamički režim}
\gusttekst\setlength{\abovedisplayskip}{2pt}\setlength{\belowdisplayskip}{2pt}\setlength{\abovedisplayshortskip}{2pt}\setlength{\belowdisplayshortskip}{2pt}\begin{columns}[c,onlytextwidth]\column{.37\textwidth}\centering\input{slike/tikz/s066_ttl_sa_izlaznom_kapacitivnoscu.tex}\column{.24\textwidth}\centering\begin{tikzpicture}[x=.6cm,y=.65cm,font=\sitneoznake,>=Latex]\draw[->](0,0)--(3.8,0)node[below]{t};\draw[->](.2,-.2)--(.2,2.1)node[left]{$v_i(t)$};\draw(.2,1.25)--(.9,1.25)--(.9,.15)--(2.7,.15)--(2.7,1.25)--(3.4,1.25);\draw[dashed](.2,.15)--(.9,.15);\node[left]at(.2,1.25){$V_H$};\node[left]at(.2,.35){$V_L$};\draw(.9,0)--(.9,-.1)node[below]{$t_1$};\draw(2.7,0)--(2.7,-.1)node[below]{$t_2$};\end{tikzpicture}
\par\medskip\begin{tikzpicture}[x=.65cm,y=.65cm,font=\sitneoznake,>=Latex]\ctikzset{bipoles/length=.5cm}\node[draw,circle,minimum size=4mm,inner sep=0pt](s)at(0,0){};\draw[->](s.north)--(s.south);\draw(s.north)--(0,.8)--(2,.8)node[ocirc]{}node[right]{$V_o(t)$};\draw(s.south)--(0,-1)node[ground]{};\draw(1.5,.8)to[C,l=$C_O$](1.5,-1)node[ground]{};\node[left]at(-.35,0){$\beta_F I_{B4}$};\end{tikzpicture}
\column{.37\textwidth}
Sa datim brojnim vrednostima tranzistor će raditi u zasićenju kada se dostigne 50\% promene pa možemo odmah da pišemo
\[t_{pLH}=\tau_2\ln\left(\frac{v(t_1^+)-v(\infty)}{\dfrac{V_{OL}+V_{OH}}2-v(\infty)}\right)\]
\end{columns}
\begin{columns}[c,onlytextwidth]\column{.44\textwidth}\centering\begin{tikzpicture}[x=.6cm,y=.24cm,font=\sitneoznake,>=Latex]
\draw[->](0,4.8)--(5,4.8)node[below]{t};\draw[->](.2,4.5)--(.2,7.0)node[above]{$v_i(t)$};\draw(.2,6.5)--(1,6.5)--(1,5.1)--(2.8,5.1)--(2.8,6.5)--(4.7,6.5);\node[left]at(.2,6.5){$V_H$};\node[left]at(.2,5.1){$V_L$};
\draw[->](0,.7)--(5,.7)node[below]{t};\draw[->](.2,.4)--(.2,3.6);\node at(1.8,3.6){$v_o(t)$};\foreach\y/\l in{2.8/V_{OH},1/V_{OL}}{\draw[dashed](.2,\y)--(3.3,\y);\node[left]at(.2,\y){$\l$};}
\draw(.2,1)--(1,1)..controls(1.2,2.5)and(1.7,2.8)..(2.2,2.8)--(2.8,2.8)--(3.05,1)--(4.7,1);
\foreach\x/\l in{1/1,2.8/2}{\draw[dashed](\x,.45)--(\x,5);\node[below]at(\x,.45){$t_\l$};}\draw[->](4.5,2.3)--(2.95,1.8);\node[right,align=left]at(4.5,2.3){Linearna\\zavisnost};\draw[->](1.08,1.5)--(1.45,-1);\node[below,align=center]at(.5,-1){Totem pol radi\\u režimu\\zasićenja};\draw[->](1.8,2.7)--(3.1,-1);\node[below,align=center]at(4.2,-1){Totem pol radi\\u aktivnom\\režimu};\draw[dashed](.2,1.9)--(3.3,1.9);\node[fill=white,inner sep=1pt]at(2,1.9){$50\%$};\end{tikzpicture}

\column{.54\textwidth}
\[v_o(t)=v_o(t_2^+)-\frac{I_{OL}}C(t-t_2)\quad za\ t\geq t_2\]
\[\frac{V_{OL}+V_{OH}}2=V_{OH}-\frac{I_{OL}}C t_{pHL}\]
\[t_{pHL}=C\frac{V_{OH}-V_{OL}}{2I_{OL}}\]
\end{columns}
\end{frame}
